\documentclass[11pt,a4paper]{article}
\usepackage{fontspec}
\setmainfont{Noto Serif CJK JP}
\newfontfamily\jpsans{Noto Sans CJK JP}
\XeTeXlinebreaklocale "ja"
\XeTeXlinebreakskip = 0pt plus 1pt minus 0.1pt
\XeTeXlinebreakpenalty = 0
\usepackage{amsmath,amssymb}
\usepackage[margin=25mm]{geometry}
\usepackage{array}
\usepackage{booktabs}
\usepackage{enumitem}
\usepackage[hidelinks]{hyperref}
\usepackage{titlesec}
\titleformat{\section}{\large\bfseries\jpsans}{\thesection}{0.8em}{}
\titleformat{\subsection}{\normalsize\bfseries\jpsans}{\thesubsection}{0.6em}{}
\setlength{\parskip}{4pt}
\renewcommand{\baselinestretch}{1.08}
\newcommand{\eP}{\ell_P}

\title{\bfseries プランク球宇宙論（Planck Sphere Cosmology）：\\プランク長スケールの球状状態を中心イメージとした統合的再提示と、\\その「言えること・言えないこと」}
\author{石井 努（Tsutomu Ishii）\\ Independent Researcher, Nagano, Japan\\ ORCID: 0009-0001-3019-3929}
\date{2026年6月23日}

\begin{document}
\maketitle

\begin{abstract}
本稿は、宇宙の最小単位を半径 $\eP$（プランク長）の球とみなし、その表面積を媒介とする核心等式 $A_P c^3 = 4\pi\hbar G$ を出発点として現代物理を統一的に再整理しようとするプランク球宇宙論（Planck Sphere Cosmology, 以下 PSC）について、その内容がどの水準で成立するかを著者自身が批判的に検討したものである。各関係式を「恒等変形（L1）／既存式への代入（L2）／次元解析（L3）」の三水準に分類し、高精度数値計算（mpmath 40桁）と記号計算（sympy）で裏づけた。結果として、PSC が厳密な意味で導出するのはプランク単位系の代数的書き換え（L1）に限られること、宇宙定数やダークセクターは単一スケールの対象であるプランク球には含まれないこと、そして完全な球ではワイルテンソルが恒等的に零となりペンローズの低エントロピー初期条件と厳密に一致すること（ただし仮説の幾何的再記述）を確認する。さらに本稿は、観測されるバリオン光子比をよく再現するパラメータゼロの経験的関係式 $\eta = \alpha^3/(8\pi)^2$（電磁結合の3乗を重力の規格化定数の2乗で割った数）を取り上げ、その分母の起源が確定する一方で生成機構は未導出であること、既知の枠組みからは自然に導けないこと、そしてその身分（偶然・構造的必然・未知の一般原理の現れ）が現時点で未確定であることを誠実に整理する。最後に PSC を無境界仮説・ワイル曲率仮説・ループ量子宇宙論の中に位置づける。本稿は誇張を排し、言えることと言えないことを明示することを第一義とする。
\end{abstract}

\section{序論：動機と中心イメージ}
PSC は、宇宙の最小単位を半径 $\eP$ の球（プランク球）とし、その表面積 $A_P = 4\pi\eP^2$ を媒介として
\begin{equation}
A_P \times c^3 = 4\pi \times \hbar \times G
\label{eq:core}
\end{equation}
が成り立つことを核心等式に据える。本稿の目的は、ここから現代物理が「導出できる」という主張が実際にどの水準で成立するのかを批判的に画定し、あわせて PSC に付随する一つの際立った数値的関係式を誠実に位置づけることである。物理理論の価値は、何を主張できるかと同程度に、何を主張できないかを正確に画定できるかにかかっている。そこで本稿では一貫して次の三水準を用いる。

\begin{description}[leftmargin=2.2em,style=nextline]
\item[L1（恒等変形）] 核心等式や定義からの代数的な書き換え。数学的に厳密だが新規の物理内容を含まない。
\item[L2（既存式への代入）] 既存物理の関係式に PSC の量を代入したもの。PSC 独自の導出ではない。
\item[L3（次元解析＋係数の後付け）] 比例関係やオーダーは出るが係数は観測から決める。厳密には導出でなくスケーリング則の提案にとどまる。
\end{description}

最大の危険は L3 を L1 と取り違えることである。以下では各関係式がどの水準かを明示し、高精度数値計算と記号計算で確認する。立場としては、潰せるものは潰し、伸ばせるものだけを伸ばす。

\paragraph{中心となる作業仮説。}本稿が中心イメージとして置くのは、\emph{宇宙初期に半径プランク長程度の球状領域が存在した可能性}である。この球状状態の内部・表面・あるいは外部への作用として、後の宇宙を特徴づける基本量（エネルギー・電荷・曲率・情報など）が潜在的に含まれていた、という発想を出発点とする。これはあくまで探求のための作業仮説であり、物理法則の主張ではない。重要なのは次の三点を初めから固定しないことである。すなわち、(i) 球の内部が真空であると決めつけない（ブラックホール内部・初期宇宙・プランク時代について確立した描像は存在せず、半径 $\eP$ の球の内部が真空でなければならない理由はない）、(ii) 量の配置をめぐる三層構造（内部・表面・外部への作用に体積因子 $4\pi/3$・表面因子 $4\pi$・規格化因子 $8\pi$ を割り当てる整理）は法則ではなく便宜的な整理の枠とみなす、(iii) したがって本質的な問いは「球が存在したか」よりも「\emph{その中に何が潜在していたか}」（情報・場・曲率・低エントロピー）である。本稿はこの作業仮説のもとで「何が厳密に言え、何が言えないか」を画定する。

\section{硬い核：核心等式とプランク量の代数（L1）}
まず核心等式 \eqref{eq:core} そのものの身分を問う。プランク長の標準的定義は $\eP^2 = \hbar G/c^3$ である。これと $A_P = 4\pi\eP^2$ を式 \eqref{eq:core} の左辺に代入すると
\begin{equation}
A_P c^3 = 4\pi\eP^2 c^3 = 4\pi\!\left(\frac{\hbar G}{c^3}\right)\! c^3 = 4\pi\hbar G
\end{equation}
となり、右辺と恒等的に一致する。数値的にも両辺は40桁の精度で一致する（相対差ゼロ）。すなわち核心等式は独立した物理的仮説ではなく、プランク長の定義の幾何的書き換えにすぎない（L1）。

係数 $4\pi$ は球の立体角であり、「球であると決めた」ことによって導入される。定義 $\eP^2=\hbar G/c^3$ 自体には幾何は含まれていない。したがって $4\pi$ という幾何因子は予言を生まない。

同様に、$\hbar = E_P t_P$、$q_P^2 = 4\pi\varepsilon_0\hbar c$、$E_P = k_B T_P$、$e = q_P\sqrt{\alpha}$ といった等価関係はいずれも代入すると $4\pi\hbar G$ に戻る恒等式であり、数値検算でも相対差はゼロであった。とくに本稿で後に用いる二つの無次元量
\begin{equation}
\alpha = \left(\frac{e}{q_P}\right)^2 \;(\text{電磁結合定数}),\qquad
8\pi = \left(\frac{m_P}{M_{\rm red}}\right)^2 \;(\text{重力の規格化定数})
\label{eq:alpha8pi}
\end{equation}
は厳密な代数関係である。ここで $q_P$ はプランク電荷、$M_{\rm red}=m_P/\sqrt{8\pi}$ は換算プランク質量であり、$8\pi$ はアインシュタイン方程式に現れる $8\pi G$ の $8\pi$、すなわち一般相対論における重力の規格化そのものである。

この監査から、PSC の関係式は構造的に三層に分かれる（表~\ref{tab:audit}）。第一に\textbf{L1 群}（核心等式・基本定数・プランク量・各等価式）は、すべてプランク長定義の恒等変形であり、既知のプランク単位系の再記述である。第二に\textbf{L2 群}（情報量・ダークマター密度）は既存物理への代入である。第三に\textbf{L3 群}（宇宙定数・ダークエネルギー密度）は次元解析にとどまる。したがって「プランク球から現代物理のほぼ全てが導出できる」という主張は、厳密な意味では成立しない。厳密な導出と呼べるのはプランク単位系の代数的書き換えのみである。これは理論を否定するものではなく、その主張範囲を防御可能な核へ正しく画定する作業である。

\begin{table}[t]
\centering
\small
\begin{tabular}{@{}lll@{}}
\toprule
関係式 & 内容 & 水準 \\
\midrule
核心等式 & $A_P c^3 = 4\pi\hbar G$ & L1（$\eP^2=\hbar G/c^3$ の書き換え）\\
基本定数 & $c,\,G,\,\hbar$ を解く & L1（代数変形）\\
プランク量 & $\eP,\,t_P,\,E_P,\,m_P$ & L1 \\
電磁気等価 & $A_P c^3 = q_P^2 G/\varepsilon_0 c$ & L1（$q_P$ の定義の代入）\\
熱統計等価 & $A_P c^3 = 4\pi k_B T_P t_P G$ & L1（$E_P=k_B T_P$ の代入）\\
情報量 & $I = A_P/4\eP^2 = \pi$ & L2（幾何比＋ベッケンシュタイン）\\
力の分離 & $A_P c^3 = e^2 G/\alpha\varepsilon_0 c$ & L1（分離は標準場の理論）\\
宇宙定数 & $\Lambda \sim 1/R^2$ & L3（次元解析）\\
ダークマター密度 & $\rho_{\rm DM}\sim H^2/4\pi G$ & L2（臨界密度の係数違い）\\
ダークエネルギー密度 & $\rho_{\rm DE}\sim 2\pi\hbar/A_P^2 c$ & L3（プランク密度級・スケール問題）\\
\bottomrule
\end{tabular}
\caption{PSC の主要関係式の導出水準。}
\label{tab:audit}
\end{table}

\section{プランク球に「含まれていないもの」：宇宙定数とダークセクター}
半径 $\eP$ のプランク球が内在する長さスケールは $\eP$ ただ一つである。一方、宇宙定数やダークエネルギーは第二のスケール、すなわち宇宙論的スケール $R_H$（ハッブル半径）の量である。両者の比は
\begin{equation}
\frac{R_H}{\eP}\approx 8.5\times10^{60},\qquad \left(\frac{R_H}{\eP}\right)^2\approx 10^{122}
\end{equation}
であり、いわゆる「宇宙定数問題の122桁」はこの比の二乗そのものである。単一スケール $\eP$ のプランク球は、第二スケールも、それらをつなぐ階層比も生み出せない。したがって宇宙定数とダークセクターは、第二スケールを生む動力学を欠く静的なプランク球には含まれない。欠けているのは、プランクスケールと宇宙論スケールをつなぐ時間発展の動力学である。

\section{情報・低エントロピーと球の幾何}
一つの明確な幾何学的対応がある。閉じたフリードマン計量（空間が3次元球）では、ワイルテンソルの全成分が恒等的に零となる（sympy による記号計算で確認）。ワイルテンソルは物質分布で決まるリッチテンソルと対照的に「自由な重力の自由度」を表し、ペンローズはこれを重力エントロピーに対応づけ、初期特異点でワイルが零であること（＝低エントロピー初期条件）を要請した（ワイル曲率仮説）。したがって「完全な球 $\Leftrightarrow$ ワイル零 $\Leftrightarrow$ 重力エントロピー最小」は厳密に正しい対応である。

ただしこれは新たな導出ではない。PSC は球を作業仮説として置いているにすぎず、低エントロピーの\emph{理由}を導いてはいない。これはワイル曲率仮説の幾何的再記述（L2）にとどまる。なお情報量 $I=A_P/4\eP^2=\pi$ は、半径 $\eP$・質量 $m_P$ の最小ブラックホールのエントロピー $S=A/4\eP^2=\pi$ と一致する。

\section{宇宙初期像：確立した枠組みの中の位置づけ}
PSC が描く「実在する小さな低エントロピーの初期状態」という像は、独立した三つの確立した枠組みのいずれとも近い。

\textbf{ペンローズのワイル曲率仮説・共形サイクリック宇宙論。}初期特異点でワイル零（共形平坦）＝低エントロピー。PSC の「球＝滑らか＝ワイル零＝低エントロピー」はこれと同一である。

\textbf{ハートル–ホーキングの無境界仮説。}宇宙の波動関数を経路積分で与え、ビッグバン特異点を滑らかな幾何で置き換える。その原型インスタントンはまさにユークリッド4次元球であり、低エントロピーの初期状態が基底状態から自然に出る。ただしこの球は虚時間（ユークリッド符号）の鞍点であり、トンネル確率を計算する数学的道具であって、実在のローレンツ的対象ではない。

\textbf{ループ量子宇宙論。}特異点を量子バウンスで解消し、エネルギー密度が約 $0.41\,\rho_{\rm Planck}$ に達するとバウンスが起こる。ただしバウンス時の体積は物質量に比例し、$\eP$ には固定されない。

以上から、PSC はこれら三つの交差点に位置する。プランク球が初期状態として一意・必然であることは、いずれの枠組みでも導出できない。PSC が単独でコミットしているのは「\emph{実在する・ローレンツ符号の・半径がちょうど $\eP$ に固定された球を、文字通りの初期状態とする}」という一点だが、これは現状では動機のない仮定である。なお、無境界仮説（ユークリッド的・滑らかな基底状態）とループ量子宇宙論（ローレンツ的・実在のバウンス）については、Bojowald と Brahma の研究が、両者を接続する一つの可能性（プランク領域での非特異な動的符号変化）を示している。

\section{際立った経験的関係式：電磁の3乗を重力の2乗で割った数}
本節では、PSC に付随する一つの数値的関係式を取り上げ、その身分を誠実に画定する。式 \eqref{eq:alpha8pi} の二つの無次元量、電磁結合 $\alpha$ と重力の規格化 $8\pi$ を用いると、次の量が定義できる。
\begin{equation}
\eta \;=\; \frac{\alpha^3}{(8\pi)^2} \;=\; \frac{\alpha^3}{64\pi^2}
\;=\; \left(\frac{e}{q_P}\right)^{6}\!\left(\frac{M_{\rm red}}{m_P}\right)^{4}.
\label{eq:eta}
\end{equation}
これは文字通り「\emph{電磁結合の3乗を、重力の規格化定数の2乗で割った数}」である。低エネルギーでの電磁結合 $\alpha=\alpha(0)=1/137.036$ を入れると
\begin{equation}
\eta = 6.152\times10^{-10}
\end{equation}
となり、観測される宇宙のバリオン光子比 $n_b/n_\gamma \approx 6.12\times10^{-10}$ と約 $0.5\%$ で一致する。式 \eqref{eq:eta} は\textbf{自由なパラメータを一切含まず}、その意味で反証可能である。

\subsection{確定している部分：分母の起源}
分母 $64\pi^2$ の起源は確定している。式 \eqref{eq:alpha8pi} より
\begin{equation}
64\pi^2 = (8\pi)^2 = \left(\frac{m_P}{M_{\rm red}}\right)^{4}
\end{equation}
であり、これは換算プランク質量の定義から一意に定まる厳密な代数関係である。すなわち分母は一般相対論の $8\pi$ の二乗であって、後から合わせた数ではない。

\subsection{確定していない部分：機構は未導出}
一方で、なぜ「電磁の3乗」と「重力の2乗」がこの形で組み合わさって観測値を与えるのか、その物理的機構は導出されていない。我々は機構を求めて複数の方向を系統的に検討したが、いずれも式 \eqref{eq:eta} を\emph{あらかじめ答えとして入力することなしには}再現できなかった。

\begin{itemize}[leftmargin=1.4em]
\item \textbf{標準的なバリオン数生成。} 既知のバリオン数生成（大統一崩壊・電弱・レプトン数生成・曲率と結合する重力的バリオン数生成）では、観測値の小ささは CP 対称性の破れ・非平衡度・脱結合温度などの比から生じる。これらは「結合定数の3乗 $\times$ 純粋な幾何因子」という形にはならない。小ささの\emph{源泉の型}が異なるため、これらから式 \eqref{eq:eta} は自然には出ない。
\item \textbf{スペクトル幾何（熱核展開）。} 4次元のラプラシアン／ディラック作用素の熱核は、前因子として $1/(4\pi)^2 = 1/(16\pi^2)$ を自然に与える。しかしこれは PSC が必要とする $1/(64\pi^2)$ と4倍ずれており、その差を埋める $1/4$ には独立な動機がない。さらに電磁結合の3乗は純粋な幾何からは決して現れない。
\item \textbf{指数定理。} 指数定理が与える量は整数（特性数）であり、$64\pi^2$（無理数）を直接与えることはできない。
\end{itemize}
要するに、$\alpha^3$ も $1/(8\pi)^2$ も部品としては既知物理に存在するが、両者を「ちょうど積として、他に余分な因子を伴わずに」生み出す単一の既知過程は見当たらない。

\subsection{数値的一致をどう評価するか}
パラメータゼロの一致は印象的だが、過大評価を避けるため、$\alpha$ と $\pi$ から作られる単純な式の中で観測値に近いものがどれくらいあるかを数えた。許容を観測誤差程度（約 $1\%$）にとると、同程度に単純な式が数個この近傍に存在し、しかもそのうちには式 \eqref{eq:eta} よりも数値的に近いものもある。したがって\emph{数値的近さそのものは決定的ではない}。式 \eqref{eq:eta} の特長は数値ではなく、その分母が後付けでなく厳密な代数関係として強制される点、すなわち\emph{構造}にある。

\subsection{この関係式の特異性：階層を持たないこと}
もう一つの観察がある。多くの宇宙論的・素粒子的な小さな量は、巨大な階層比に支配される。たとえば宇宙定数は $(\eP/R_H)^2\sim10^{-122}$ という宇宙論的階層、ベックが提案した宇宙定数の関係式は電子コンプトン長とプランク長の比 $(m_e/m_P)^6$ という質量階層に支配される。これに対し式 \eqref{eq:eta} は\emph{いかなる階層比も含まず}、低エネルギーの結合定数の冪と純粋な幾何因子だけで小さい。この意味で $\eta$ は孤立しており、他の観測量と共通の構造（たとえば $\alpha$ の冪の族）を形成しない。ただし孤立それ自体は「深い意味」の証拠ではなく、偶然の可能性とも両立する。

\subsection{身分：三つの可能性のあいだで保留}
以上を踏まえ、式 \eqref{eq:eta} の身分は次の三つの可能性のいずれとも確定できず、本稿では\textbf{保留}とする。
\begin{enumerate}[leftmargin=1.6em]
\item \textbf{偶然} —— 単なる数値的一致（上記のとおり一致は軽度にしか稀でない）。
\item \textbf{構造的必然} —— プランク球の作業仮説に内在する構造比（現時点で積極的証拠はない）。
\item \textbf{未知の一般原理の現れ} —— まだ知られていない一般原理によって「結合定数の3乗 $\times$ 微小な幾何因子」という型が必然化される可能性。これは特定の作業仮説に固有でなく、より大きな理論の一断面でありうる。
\end{enumerate}
この保留を将来動かしうるのは、観測されるバリオン光子比の精度向上（現在の中心値が保たれたまま精度が $0.1\%$ 級に達すれば、式 \eqref{eq:eta} の $0.5\%$ のずれは検証あるいは反証の対象となる）、あるいは上記の不可能性をより厳密に示すか、逆に機構を発見することである。

\section{方法論：誠実さの規律}
本研究は一貫して一つの規律のもとで進めた。すなわち、確立した物理・定義や恒等式・数値的一致・思弁的仮説を常に区別し、未導出のものは未導出と明記することである。とくに「先に魅力的な数値一致を見つけ、後からその部品（たとえば $64$ や $\pi$ の冪）に意味を付ける」という手続きは、いくらでも後付けが効くため、それ自体を成果とはみなさない方針をとった。式 \eqref{eq:eta} の機構を探る際にも、答えを模型にあらかじめ入力しないという自己規律を課し、自然に現れたときにのみ成果と認めることとした。本稿の数値はすべて mpmath による40桁精度計算と sympy による記号計算で裏づけ、可能な限り一次資料を確認した。検討は二つの独立した解析支援（異なる大規模言語モデル）と著者の相互批判を通じて行い、結論はいずれも一致した。

\section{総括：言えること／言えないこと}
\begin{center}
\small
\begin{tabular}{@{}p{0.46\linewidth}p{0.46\linewidth}@{}}
\toprule
\textbf{言えること（確立・確定）} & \textbf{言えないこと（未導出・思弁）} \\
\midrule
核心等式・各等価式・プランク量はプランク長定義の恒等変形として厳密（L1） & それらが「新しい物理」であること（既知のプランク単位系の再記述）\\
完全な球でワイル零＝低エントロピーは厳密な対応 & 低エントロピーの\emph{理由}（球は作業仮説として置いているだけ）\\
$\eta=\alpha^3/(8\pi)^2$ が観測バリオン光子比と $0.5\%$ で一致、パラメータゼロで反証可能 & その生成機構（既知の枠組みからは自然に導けない）\\
分母 $64\pi^2=(8\pi)^2$ の起源は確定（重力の規格化の二乗） & $\eta$ が偶然か・構造的必然か・未知の一般原理の現れか（保留）\\
宇宙定数・ダークセクターは単一スケールの球には含まれない & 第二スケールへの時間発展の動力学 \\
PSC は三つの確立した枠組みの交差点に位置する & プランク球の一意性・実在性、独自の反証可能な予言 \\
\bottomrule
\end{tabular}
\end{center}

残るオープンな問いは、(i) 式 \eqref{eq:eta} の身分、(ii) プランク球の内部に何が潜在していたか（情報・場・曲率・低エントロピー）、(iii) なぜ初期状態が半径 $\eP$ の実在する球でなければならないかの動的根拠、(iv) プランクスケールから宇宙論スケールへの展開則、である。これらはいずれも現代物理の未解決領域に属する。

\section{結論}
本稿はプランク球宇宙論を著者自身が批判的に検討した。第一に、PSC が厳密に導出するのはプランク単位系の代数的書き換え（L1）に限られる。第二に、宇宙定数とダークセクターは静的な単一スケールのプランク球には含まれない。第三に、完全な球でワイルテンソルが零となり低エントロピー初期条件と一致するという積極的事実は成立するが、これはワイル曲率仮説の幾何的再記述である。第四に、PSC はワイル曲率仮説・無境界仮説・ループ量子宇宙論の交差点に位置する統合的再提示として正直に提示しうる。そして第五に、電磁の3乗を重力の2乗で割った数 $\eta=\alpha^3/(8\pi)^2$ は観測バリオン光子比をパラメータなしでよく再現する経験的関係式であり、分母の起源は確定するが、その機構は未導出で、身分は保留である。本稿は誇張を排し、プランク球という一つのイメージのもとで「ここまでは言える・ここから先はまだ言えない」を画定することを目的とした。統一は、導出としてではなく、複数の確立した描像と一つの際立った経験則を見渡す\emph{中心イメージ}として提示される。

\section*{付録：方法と検算、棄却の記録}
\textbf{方法と検算。}数値はすべて Python の mpmath により40桁精度で計算し、CODATA 値を用いた。ワイルテンソルは sympy による記号計算で閉じたフリードマン計量の全成分が零であることを確認した。核心等式・各等価式の恒等性、$64\pi^2=(8\pi)^2=(m_P/M_{\rm red})^4$、4次元熱核の前因子 $1/(16\pi^2)$、$\alpha$--$\pi$ 単純式の探索などをこれらの計算で裏づけた。

\textbf{棄却の記録。}検討の過程で次は棄却または保留とした。重力的バリオン数生成がそのまま式 \eqref{eq:eta} を導くという仮説（小ささの源泉が温度比であり $\alpha^3$ でない）、$\eta$ の分子の $3$ を世代数とみなす解釈、分母の $16$ をスピノル構造とみなす解釈、$\eta=\alpha(\alpha/8\pi)^2$ という分解（形はアノマリー係数に似るが単一過程に対応せず、後付けにとどまる）。なお、実在・半径 $\eP$ 固定の球にコミットするとバウンス時の密度が $3/(4\pi)\approx0.239\,\rho_{\rm Planck}$ と定まり、ループ量子宇宙論の $0.41\,\rho_{\rm Planck}$ とは異なる数になるが、これは観測量ではなく仮定にも依存するため、確立した予言ではなく\emph{探求の芽}にとどまる。

\begin{thebibliography}{9}
\small
\bibitem{HH} J. B. Hartle, S. W. Hawking, \textit{Wave function of the Universe}; レビュー: \textit{Review of the No-Boundary Wave Function}, arXiv:2303.08802.
\bibitem{saddle} J. Feldbrugge, J.-L. Lehners ほか, \textit{Path integrals, saddle points and the beginning of the universe}, arXiv:2303.11450（無境界鞍点がユークリッド4次元球 $S^4$）.
\bibitem{Penrose} R. Penrose, ワイル曲率仮説・共形サイクリック宇宙論; \textit{CCC, gravitational entropy and quantum information}, Gen. Rel. Grav. (2023).
\bibitem{LQC} A. Ashtekar ほか, \textit{Singularity Resolution in Loop Quantum Cosmology: A Brief Overview}, arXiv:0812.4703.
\bibitem{LQCvol} \textit{On the evolution of the volume in Loop Quantum Cosmology}, arXiv:2311.18066（バウンス体積の物質依存性）.
\bibitem{BB} M. Bojowald, S. Brahma, \textit{Loops Rescue the No-Boundary Proposal}, Phys. Rev. Lett. \textbf{121}, 201301 (2018), arXiv:1810.09871.
\bibitem{Beck} C. Beck, \textit{Cosmological constant, fine structure constant and beyond}, Eur. Phys. J. C (2016), arXiv:1605.04571.
\bibitem{baryo} レビュー: \textit{Baryogenesis} (バリオン数生成の総説), Rev. Mod. Phys. \textbf{93}, 035004 (2021).
\end{thebibliography}

\end{document}
